\documentclass[aspectratio=169,12pt]{beamer}

% ============================================
% CONFIGURAÇÕES DO TEMA
% ============================================
\usetheme{Madrid}
\usecolortheme{whale}

% Cores personalizadas (azul institucional)
\definecolor{unespazul}{RGB}{0,61,121}
\definecolor{unespvermelho}{RGB}{178,34,34}
\setbeamercolor{structure}{fg=unespazul}
\setbeamercolor{title}{fg=white,bg=unespazul}
\setbeamercolor{frametitle}{fg=white,bg=unespazul}
\setbeamercolor{block title}{fg=white,bg=unespazul}
\setbeamercolor{block body}{bg=unespazul!10}

% ============================================
% PACOTES
% ============================================
\usepackage[utf8]{inputenc}
\usepackage[T1]{fontenc}
\usepackage[brazil]{babel}
\usepackage{amsmath,amssymb,amsfonts}
\usepackage{graphicx}
\usepackage{tikz}
\usetikzlibrary{shapes,arrows,positioning,calc,decorations.pathreplacing}
\usepackage{booktabs}
\usepackage{multimedia}
\usepackage{hyperref}

% ============================================
% CONFIGURAÇÕES ADICIONAIS
% ============================================
\setbeamertemplate{navigation symbols}{}
\setbeamertemplate{footline}[frame number]
\setbeamertemplate{caption}[numbered]

% Numerar seções no TOC
\setbeamertemplate{section in toc}[sections numbered]
\setbeamertemplate{subsection in toc}[subsections numbered]

% ============================================
% INFORMAÇÕES DA APRESENTAÇÃO
% ============================================
\title[Mecânica Clássica e IA]{A Relevância da Mecânica Clássica na Era da Inteligência Artificial Generativa}
\subtitle{Aula Magna -- Concurso para Professor Titular}
\author{Candidato}
\institute[UNESP]{
    Departamento de Biofísica e Farmacologia\\
    Instituto de Biociências\\
    UNESP Botucatu
}
\date{\today}

% Logo (descomentar se tiver o arquivo)
% \titlegraphic{\includegraphics[height=1.5cm]{unesp-logo.png}}

\begin{document}

% ============================================
% SLIDE TÍTULO
% ============================================
\begin{frame}
    \titlepage
\end{frame}

% ============================================
% EPÍGRAFE
% ============================================
\begin{frame}
    \centering
    \vfill
    \begin{quote}
        \Large\itshape
        ``Se vi mais longe, foi por estar sobre ombros de gigantes.''
    \end{quote}
    \vspace{0.5cm}
    \hfill --- Isaac Newton, 1675
    \vfill
\end{frame}

% ============================================
% SUMÁRIO
% ============================================
\begin{frame}{Roteiro}
    \tableofcontents
\end{frame}

% ============================================
% SEÇÃO 1: INTRODUÇÃO FILOSÓFICA
% ============================================
\section{Introdução Filosófica}

\begin{frame}{A Questão Epistemológica Fundamental}
    \begin{center}
        \Large\textbf{O que significa ``compreender'' a natureza?}
    \end{center}
    \vspace{1cm}
    \begin{columns}[T]
        \begin{column}{0.48\textwidth}
            \begin{block}{Compreensão Analítica}
                \begin{itemize}
                    \item Leis fundamentais
                    \item Relações causais
                    \item Dedução matemática
                    \item \textbf{Por que} acontece
                \end{itemize}
            \end{block}
        \end{column}
        \begin{column}{0.48\textwidth}
            \begin{block}{Compreensão Estatística}
                \begin{itemize}
                    \item Padrões nos dados
                    \item Correlações
                    \item Inferência probabilística
                    \item \textbf{O que} acontece
                \end{itemize}
            \end{block}
        \end{column}
    \end{columns}
\end{frame}

\begin{frame}{Modelos Analíticos vs. Modelos Estatísticos}
    \begin{center}
    \begin{tikzpicture}[scale=0.9, transform shape,
        box/.style={rectangle, draw=unespazul, fill=unespazul!20, thick,
                    minimum width=3cm, minimum height=1.2cm, align=center}]

        % Lado esquerdo - Analítico
        \node[box] (leis) at (-4,2) {Leis\\Fundamentais};
        \node[box] (modelo) at (-4,0) {Modelo\\Matemático};
        \node[box] (pred1) at (-4,-2) {Previsões};

        \draw[->, thick, unespazul] (leis) -- (modelo);
        \draw[->, thick, unespazul] (modelo) -- (pred1);

        \node[above=0.3cm of leis, font=\bfseries] {Abordagem Analítica};

        % Lado direito - Estatístico
        \node[box] (dados) at (4,2) {Dados\\Massivos};
        \node[box] (padroes) at (4,0) {Padrões\\Aprendidos};
        \node[box] (pred2) at (4,-2) {Previsões};

        \draw[->, thick, unespazul] (dados) -- (padroes);
        \draw[->, thick, unespazul] (padroes) -- (pred2);

        \node[above=0.3cm of dados, font=\bfseries] {Abordagem Estatística};

        % Setas de comparação
        \draw[<->, dashed, thick, unespvermelho] (-1.5,0) -- (1.5,0)
            node[midway, above, font=\small] {Complementares?};
    \end{tikzpicture}
    \end{center}
\end{frame}

\begin{frame}{Duas Tradições de Conhecimento}
    \begin{columns}[T]
        \begin{column}{0.48\textwidth}
            \begin{block}{Tradição Analítica (Top-Down)}
                \begin{itemize}
                    \item Princípios primeiros
                    \item Dedução rigorosa
                    \item Leis universais
                    \item Determinismo
                \end{itemize}
                \vspace{0.3cm}
                \centering
                \textbf{Mecânica Clássica}
            \end{block}
        \end{column}
        \begin{column}{0.48\textwidth}
            \begin{block}{Tradição Empírica (Bottom-Up)}
                \begin{itemize}
                    \item Dados observados
                    \item Indução estatística
                    \item Padrões locais
                    \item Probabilismo
                \end{itemize}
                \vspace{0.3cm}
                \centering
                \textbf{IA Generativa}
            \end{block}
        \end{column}
    \end{columns}
    \vspace{0.5cm}
    \begin{center}
        \textit{Ambas buscam modelos preditivos da realidade}
    \end{center}
\end{frame}

\begin{frame}{Paralelo Histórico}
    \begin{center}
    \begin{tikzpicture}[scale=0.85, transform shape]
        % Timeline
        \draw[thick, ->, unespazul] (0,0) -- (12,0) node[right] {tempo};

        % Marcadores
        \draw[thick, unespazul] (1,0.2) -- (1,-0.2) node[below] {1687};
        \draw[thick, unespazul] (4,0.2) -- (4,-0.2) node[below] {1788};
        \draw[thick, unespazul] (6,0.2) -- (6,-0.2) node[below] {1833};
        \draw[thick, unespazul] (9,0.2) -- (9,-0.2) node[below] {2017};
        \draw[thick, unespazul] (11,0.2) -- (11,-0.2) node[below] {2022};

        % Eventos
        \node[above=0.5cm, align=center, font=\small] at (1,0) {Newton\\Principia};
        \node[above=0.5cm, align=center, font=\small] at (4,0) {Lagrange\\Méc. Anal.};
        \node[above=0.5cm, align=center, font=\small] at (6,0) {Hamilton\\Formalismo};
        \node[above=0.5cm, align=center, font=\small] at (9,0) {Transformer\\Vaswani};
        \node[above=0.5cm, align=center, font=\small] at (11,0) {ChatGPT\\OpenAI};

        % Arcos
        \draw[thick, unespvermelho, decorate, decoration={brace, amplitude=10pt}]
            (0.5,1.8) -- (6.5,1.8) node[midway, above=12pt] {Mecânica Clássica};
        \draw[thick, unespvermelho, decorate, decoration={brace, amplitude=10pt}]
            (8.5,1.8) -- (11.5,1.8) node[midway, above=12pt] {IA Generativa};
    \end{tikzpicture}
    \end{center}
    \vspace{0.5cm}
    \begin{center}
        \textit{Dois paradigmas de formalização do conhecimento sobre a natureza}
    \end{center}
\end{frame}

% ============================================
% SEÇÃO 2: FUNDAMENTOS DA MECÂNICA CLÁSSICA
% ============================================
\section{Fundamentos da Mecânica Clássica}

\begin{frame}{A Revolução Newtoniana}
    \begin{columns}[T]
        \begin{column}{0.55\textwidth}
            \textbf{As Três Leis de Newton}
            \vspace{0.3cm}

            \textbf{1ª Lei} (Inércia):
            \begin{equation*}
                \vec{F} = 0 \Rightarrow \vec{v} = \text{constante}
            \end{equation*}

            \textbf{2ª Lei} (Dinâmica):
            \begin{equation*}
                \vec{F} = m\vec{a} = m\frac{d^2\vec{r}}{dt^2}
            \end{equation*}

            \textbf{3ª Lei} (Ação e Reação):
            \begin{equation*}
                \vec{F}_{12} = -\vec{F}_{21}
            \end{equation*}
        \end{column}
        \begin{column}{0.42\textwidth}
            \begin{block}{Paradigma de Simplicidade}
                \begin{itemize}
                    \item Três leis
                    \item Poder explicativo universal
                    \item Matematização da natureza
                \end{itemize}
            \end{block}
            \vspace{0.3cm}
            \begin{alertblock}{Gravitação Universal}
                \begin{equation*}
                    F = G\frac{m_1 m_2}{r^2}
                \end{equation*}
            \end{alertblock}
        \end{column}
    \end{columns}
\end{frame}

\begin{frame}{O Determinismo Laplaciano}
    \begin{quote}
        \itshape
        ``Uma inteligência que, em um dado instante, conhecesse todas as forças que animam a natureza e a situação de cada ser que a compõe [...] nada seria incerto para ela, e o futuro, como o passado, estaria presente aos seus olhos.''
    \end{quote}
    \hfill --- Pierre-Simon Laplace, 1814

    \vspace{0.5cm}

    \begin{block}{Implicações Filosóficas}
        \begin{itemize}
            \item Conhecendo $\vec{r}(t_0)$ e $\vec{v}(t_0)$ $\Rightarrow$ conhecemos $\vec{r}(t)$ para todo $t$
            \item O universo como um relógio mecânico
            \item Previsibilidade total em princípio
        \end{itemize}
    \end{block}
\end{frame}

\begin{frame}{Formalismo Lagrangiano}
    \begin{block}{Princípio de Mínima Ação (Hamilton)}
        \begin{equation*}
            \delta S = \delta \int_{t_1}^{t_2} L(q, \dot{q}, t) \, dt = 0
        \end{equation*}
    \end{block}

    \begin{columns}[T]
        \begin{column}{0.48\textwidth}
            \textbf{Lagrangiana:}
            \begin{equation*}
                L = T - V
            \end{equation*}

            \textbf{Equações de Euler-Lagrange:}
            \begin{equation*}
                \frac{d}{dt}\frac{\partial L}{\partial \dot{q}_i} - \frac{\partial L}{\partial q_i} = 0
            \end{equation*}
        \end{column}
        \begin{column}{0.48\textwidth}
            \begin{block}{Vantagens}
                \begin{itemize}
                    \item Coordenadas generalizadas
                    \item Invariância por transformações
                    \item Elegância matemática
                    \item Conexão com simetrias
                \end{itemize}
            \end{block}
        \end{column}
    \end{columns}
\end{frame}

\begin{frame}{Formalismo Hamiltoniano}
    \begin{columns}[T]
        \begin{column}{0.48\textwidth}
            \textbf{Hamiltoniana:}
            \begin{equation*}
                H(q, p, t) = \sum_i p_i \dot{q}_i - L
            \end{equation*}

            \textbf{Equações Canônicas:}
            \begin{align*}
                \dot{q}_i &= \frac{\partial H}{\partial p_i} \\[0.3cm]
                \dot{p}_i &= -\frac{\partial H}{\partial q_i}
            \end{align*}
        \end{column}
        \begin{column}{0.48\textwidth}
            \begin{block}{Espaço de Fase}
                \begin{itemize}
                    \item Espaço $(q, p)$ de $2n$ dimensões
                    \item Trajetórias determinísticas
                    \item Fluxo hamiltoniano
                    \item Base para Mecânica Quântica
                \end{itemize}
            \end{block}

            \begin{alertblock}{Conservação}
                \begin{equation*}
                    \frac{dH}{dt} = \frac{\partial H}{\partial t}
                \end{equation*}
            \end{alertblock}
        \end{column}
    \end{columns}
\end{frame}

\begin{frame}{Problema dos Dois Corpos}
    \begin{columns}[T]
        \begin{column}{0.55\textwidth}
            \begin{block}{Solução Analítica Exata}
                Redução ao problema de um corpo:
                \begin{equation*}
                    \mu = \frac{m_1 m_2}{m_1 + m_2}
                \end{equation*}

                Órbitas cônicas:
                \begin{equation*}
                    r(\theta) = \frac{p}{1 + e\cos\theta}
                \end{equation*}
            \end{block}

            \textbf{Leis de Kepler} derivadas matematicamente
        \end{column}
        \begin{column}{0.42\textwidth}
            \begin{center}
            \begin{tikzpicture}[scale=0.7]
                % Órbita elíptica
                \draw[thick, unespazul] (0,0) ellipse (2.5cm and 1.5cm);
                % Sol
                \fill[yellow!80!orange] (-1,0) circle (0.2cm);
                \node[below] at (-1,-0.3) {$m_1$};
                % Planeta
                \fill[unespazul] (2.5,0) circle (0.15cm);
                \node[right] at (2.7,0) {$m_2$};
                % Vetor r
                \draw[->, thick, unespvermelho] (-1,0) -- (2.5,0) node[midway, above] {$\vec{r}$};
            \end{tikzpicture}
            \end{center}

            \vspace{0.3cm}
            \textit{Triunfo da mecânica newtoniana}
        \end{column}
    \end{columns}
\end{frame}

\begin{frame}{Problema dos Três Corpos}
    \begin{columns}[T]
        \begin{column}{0.55\textwidth}
            \begin{alertblock}{Sem Solução Analítica Geral}
                \begin{itemize}
                    \item Poincaré (1890): demonstrou a inexistência de solução geral
                    \item Sensibilidade às condições iniciais
                    \item Nascimento da teoria do caos
                \end{itemize}
            \end{alertblock}

            \vspace{0.3cm}

            \textbf{Caos Determinístico:}
            \begin{itemize}
                \item Equações determinísticas
                \item Comportamento imprevisível
                \item Expoentes de Lyapunov
            \end{itemize}
        \end{column}
        \begin{column}{0.42\textwidth}
            \begin{center}
            \begin{tikzpicture}[scale=0.6]
                % Três corpos
                \fill[yellow!80!orange] (0,0) circle (0.3cm);
                \fill[unespazul] (2,1.5) circle (0.2cm);
                \fill[unespvermelho] (-1,2) circle (0.15cm);

                % Vetores
                \draw[->, thick, gray] (0,0) -- (2,1.5);
                \draw[->, thick, gray] (0,0) -- (-1,2);
                \draw[->, thick, gray] (2,1.5) -- (-1,2);

                % Trajetórias caóticas
                \draw[thick, unespazul, dashed] (2,1.5)
                    to[out=120, in=0] (0.5,2.5)
                    to[out=180, in=90] (-1.5,1);
                \draw[thick, unespvermelho, dashed] (-1,2)
                    to[out=-30, in=150] (1,1)
                    to[out=-30, in=180] (2.5,0);
            \end{tikzpicture}
            \end{center}

            \vspace{0.3cm}
            \textit{Limites da previsibilidade}
        \end{column}
    \end{columns}
\end{frame}

\begin{frame}{Do Determinismo ao Caos}
    \begin{alertblock}{Resultado Fundamental}
        \textbf{Sistemas determinísticos e conservativos podem ser caóticos}
    \end{alertblock}
    \vspace{-0.2cm}
    \begin{columns}[T]
        \begin{column}{0.55\textwidth}
            \begin{block}{Sensibilidade às Condições Iniciais}
                Pequenas diferenças $\Rightarrow$ grandes divergências
                \begin{equation*}
                    |\delta \vec{r}(t)| \sim |\delta \vec{r}(0)| e^{\lambda t}
                \end{equation*}
                onde $\lambda > 0$ é o expoente de Lyapunov
            \end{block}
        \end{column}
        \begin{column}{0.42\textwidth}
            \begin{block}{Teorema KAM}
                \small
                Kolmogorov-Arnold-Moser
                \begin{itemize}
                    \item Estabilidade parcial
                    \item Toros invariantes
                    \item Ilhas de regularidade
                \end{itemize}
            \end{block}
        \end{column}
    \end{columns}
    \vspace{-0.1cm}
    \begin{block}{}
        \centering
        Mesmo com energia conservada, há \textbf{horizonte de previsibilidade finito}
    \end{block}
\end{frame}

% ============================================
% SEÇÃO 3: INTELIGÊNCIA ARTIFICIAL GENERATIVA
% ============================================
\section{Inteligência Artificial Generativa}

\begin{frame}{Uma Nova Forma de Conhecimento}
    \begin{center}
        \Large\textbf{Paradigma da IA: Aprendizado a partir de Padrões}
    \end{center}

    \vspace{0.5cm}

    \begin{columns}[T]
        \begin{column}{0.48\textwidth}
            \begin{block}{IA Simbólica (Clássica)}
                \begin{itemize}
                    \item Regras explícitas
                    \item Conhecimento codificado
                    \item Raciocínio dedutivo
                    \item ``GOFAI''
                \end{itemize}
            \end{block}
        \end{column}
        \begin{column}{0.48\textwidth}
            \begin{block}{IA Conexionista (Moderna)}
                \begin{itemize}
                    \item Padrões aprendidos
                    \item Conhecimento distribuído
                    \item Inferência estatística
                    \item Redes neurais profundas
                \end{itemize}
            \end{block}
        \end{column}
    \end{columns}

    \vspace{0.5cm}
    \begin{center}
        \textit{Mudança epistemológica fundamental}
    \end{center}
\end{frame}

\begin{frame}{Redes Neurais: Representação Distribuída}
    \begin{columns}[T]
        \begin{column}{0.55\textwidth}
            \textbf{Neurônio Artificial:}
            \vspace{-0.2cm}
            \begin{equation*}
                y = \sigma\left(\sum_i w_i x_i + b\right)
            \end{equation*}
            \vspace{-0.3cm}
            \textbf{Aprendizado (gradiente descendente):}
            \vspace{-0.2cm}
            \begin{equation*}
                w \leftarrow w - \eta \nabla_w \mathcal{L}
            \end{equation*}
            \vspace{-0.4cm}
            \begin{block}{Características}
                \begin{itemize}
                    \item Bilhões de parâmetros
                    \item Representações emergentes
                    \item Aproximadores universais
                \end{itemize}
            \end{block}
        \end{column}
        \begin{column}{0.42\textwidth}
            \begin{center}
            \begin{tikzpicture}[scale=0.6,
                neuron/.style={circle, draw=unespazul, fill=unespazul!30, minimum size=0.6cm}]
                % Input layer
                \foreach \i in {1,2,3} {
                    \node[neuron] (I\i) at (0, 2-\i) {};
                }
                % Hidden layer
                \foreach \i in {1,2,3,4} {
                    \node[neuron] (H\i) at (2, 2.5-\i) {};
                }
                % Output layer
                \foreach \i in {1,2} {
                    \node[neuron] (O\i) at (4, 1.5-\i) {};
                }
                % Connections
                \foreach \i in {1,2,3} {
                    \foreach \j in {1,2,3,4} {
                        \draw[->, gray] (I\i) -- (H\j);
                    }
                }
                \foreach \i in {1,2,3,4} {
                    \foreach \j in {1,2} {
                        \draw[->, gray] (H\i) -- (O\j);
                    }
                }
                % Labels
                \node[left] at (-0.5, 0.5) {\small Input};
                \node[above] at (2, 2) {\small Hidden};
                \node[right] at (4.5, 0.5) {\small Output};
            \end{tikzpicture}
            \end{center}
        \end{column}
    \end{columns}
\end{frame}

\begin{frame}{Geometrização do Conhecimento}
    \begin{center}
        \large\textbf{A Grande Ideia: Conhecimento como Geometria}
    \end{center}
    \vspace{0.2cm}
    \begin{block}{Insight Fundamental}
        Conceitos e palavras como \textbf{pontos em espaços de alta dimensão}:
        \begin{itemize}
            \item \textbf{Proximidade} $\approx$ similaridade semântica
            \item \textbf{Direções} $\approx$ relações conceituais
            \item \textbf{Operações vetoriais} $\approx$ raciocínio analógico
        \end{itemize}
    \end{block}
    \begin{columns}[T]
        \begin{column}{0.48\textwidth}
            \begin{alertblock}{Paralelo com a Física}
                Assim como a Relatividade geometrizou o espaço-tempo, a IA geometriza o \textit{espaço semântico}
            \end{alertblock}
        \end{column}
        \begin{column}{0.48\textwidth}
            \begin{block}{Implicação}
                O significado emerge da \textbf{estrutura geométrica}, não de definições explícitas
            \end{block}
        \end{column}
    \end{columns}
\end{frame}

\begin{frame}{Espaços Latentes: Geometria Oculta}
    \begin{columns}[T]
        \begin{column}{0.55\textwidth}
            \begin{block}{Representações Aprendidas}
                \begin{itemize}
                    \item Dados $\rightarrow$ espaço de alta dimensão
                    \item Estrutura geométrica significativa
                    \item Operações algébricas com sentido
                \end{itemize}
            \end{block}

            \vspace{0.3cm}

            \textbf{Exemplo (Word2Vec):}
            \begin{equation*}
                \vec{rei} - \vec{homem} + \vec{mulher} \approx \vec{rainha}
            \end{equation*}

            \vspace{0.2cm}
            \textit{Analogias como paralelogramos no espaço vetorial}
        \end{column}
        \begin{column}{0.42\textwidth}
            \begin{center}
            \begin{tikzpicture}[scale=0.7]
                % Axes
                \draw[->, thick, gray] (-0.5,0) -- (3.5,0);
                \draw[->, thick, gray] (0,-0.5) -- (0,3);

                % Points
                \fill[unespazul] (1,0.5) circle (0.1cm) node[below right] {\tiny homem};
                \fill[unespazul] (2.5,0.8) circle (0.1cm) node[below right] {\tiny rei};
                \fill[unespvermelho] (0.8,2) circle (0.1cm) node[above left] {\tiny mulher};
                \fill[unespvermelho] (2.3,2.3) circle (0.1cm) node[above right] {\tiny rainha};

                % Vectors
                \draw[->, thick, unespazul] (1,0.5) -- (2.5,0.8);
                \draw[->, thick, unespvermelho] (0.8,2) -- (2.3,2.3);
                \draw[->, dashed, gray] (1,0.5) -- (0.8,2);
                \draw[->, dashed, gray] (2.5,0.8) -- (2.3,2.3);
            \end{tikzpicture}
            \end{center}

            \vspace{0.3cm}
            \textit{Geometria com semântica}
        \end{column}
    \end{columns}
\end{frame}

\begin{frame}{Arquitetura Transformer}
    \begin{columns}[T]
        \begin{column}{0.55\textwidth}
            \textbf{Mecanismo de Atenção:}
            \begin{equation*}
                \text{Attention}(Q, K, V) = \text{softmax}\left(\frac{QK^T}{\sqrt{d_k}}\right)V
            \end{equation*}

            \begin{block}{Inovações}
                \begin{itemize}
                    \item Self-attention: contexto global
                    \item Paralelização eficiente
                    \item Escalabilidade
                    \item Transferência de aprendizado
                \end{itemize}
            \end{block}
        \end{column}
        \begin{column}{0.42\textwidth}
            \begin{alertblock}{LLMs}
                Large Language Models
                \begin{itemize}
                    \item GPT, Claude, Gemini
                    \item $10^{11}$ parâmetros
                    \item Treinados em terabytes
                    \item Capacidades emergentes
                \end{itemize}
            \end{alertblock}
        \end{column}
    \end{columns}

    \vspace{0.3cm}
    \begin{center}
        \textit{Vaswani et al., ``Attention Is All You Need'', 2017}
    \end{center}
\end{frame}

\begin{frame}{Geração vs. Compreensão}
    \begin{center}
        \Large\textbf{O Dilema da ``Caixa Preta''}
    \end{center}

    \vspace{0.5cm}

    \begin{columns}[T]
        \begin{column}{0.48\textwidth}
            \begin{block}{O que a IA faz bem}
                \begin{itemize}
                    \item Reconhecer padrões
                    \item Gerar conteúdo coerente
                    \item Interpolar entre exemplos
                    \item Generalizar estatisticamente
                \end{itemize}
            \end{block}
        \end{column}
        \begin{column}{0.48\textwidth}
            \begin{alertblock}{Limitações}
                \begin{itemize}
                    \item Não ``compreende'' causalmente
                    \item Alucinações
                    \item Raciocínio lógico limitado
                    \item Interpretabilidade restrita
                \end{itemize}
            \end{alertblock}
        \end{column}
    \end{columns}

    \vspace{0.5cm}
    \begin{center}
        \textit{``Stochastic parrots'' vs. compreensão genuína?}
    \end{center}
\end{frame}

% ============================================
% SEÇÃO 4: PONTES CONCEITUAIS
% ============================================
\section{Pontes Conceituais}

\begin{frame}{Princípios Variacionais: Conexão Profunda}
    \begin{columns}[T]
        \begin{column}{0.48\textwidth}
            \begin{block}{Mecânica Clássica}
                \textbf{Princípio de Mínima Ação:}
                \begin{equation*}
                    \delta S = \delta \int L \, dt = 0
                \end{equation*}

                A natureza ``escolhe'' o caminho que extremiza a ação
            \end{block}
        \end{column}
        \begin{column}{0.48\textwidth}
            \begin{block}{Redes Neurais}
                \textbf{Minimização da Loss:}
                \begin{equation*}
                    \min_\theta \mathcal{L}(\theta) = \min_\theta \sum_i \ell(f_\theta(x_i), y_i)
                \end{equation*}

                O modelo ``escolhe'' parâmetros que minimizam o erro
            \end{block}
        \end{column}
    \end{columns}

    \vspace{0.5cm}

    \begin{alertblock}{Princípio Unificador}
        \centering
        \textbf{Otimização} como linguagem comum
    \end{alertblock}
\end{frame}

\begin{frame}{Gradiente Descendente e Cálculo de Variações}
    \begin{columns}[T]
        \begin{column}{0.48\textwidth}
            \textbf{Cálculo de Variações:}
            \begin{equation*}
                \frac{\delta F}{\delta y} = \frac{\partial f}{\partial y} - \frac{d}{dx}\frac{\partial f}{\partial y'}
            \end{equation*}

            Encontra funções que extremizam funcionais
        \end{column}
        \begin{column}{0.48\textwidth}
            \textbf{Gradiente Descendente:}
            \begin{equation*}
                \theta_{t+1} = \theta_t - \eta \nabla_\theta \mathcal{L}
            \end{equation*}

            Encontra parâmetros que minimizam funções
        \end{column}
    \end{columns}

    \vspace{0.5cm}

    \begin{center}
    \begin{tikzpicture}[scale=0.8]
        % Surface
        \draw[thick, unespazul, domain=-2:2, samples=50]
            plot (\x, {0.5*\x*\x});

        % Gradient descent steps
        \fill[unespvermelho] (-1.8, 1.62) circle (0.08cm);
        \fill[unespvermelho] (-1.2, 0.72) circle (0.08cm);
        \fill[unespvermelho] (-0.6, 0.18) circle (0.08cm);
        \fill[unespvermelho] (0, 0) circle (0.08cm);

        \draw[->, thick, unespvermelho] (-1.8, 1.62) -- (-1.25, 0.77);
        \draw[->, thick, unespvermelho] (-1.2, 0.72) -- (-0.65, 0.23);
        \draw[->, thick, unespvermelho] (-0.6, 0.18) -- (-0.05, 0.05);

        \node[below] at (0, -0.3) {mínimo};
        \node[right] at (2.2, 2) {$\mathcal{L}(\theta)$};
    \end{tikzpicture}
    \end{center}
\end{frame}

\begin{frame}{Princípio da Energia Livre (Friston)}
    \begin{block}{Framework Unificador}
        Sistemas adaptativos minimizam \textbf{energia livre variacional}:
        \begin{equation*}
            F = D_{KL}[q(\theta) \| p(\theta | x)] - \log p(x)
        \end{equation*}
    \end{block}

    \begin{columns}[T]
        \begin{column}{0.48\textwidth}
            \textbf{Aplicações:}
            \begin{itemize}
                \item Neurociência computacional
                \item Percepção ativa
                \item Controle motor
                \item Aprendizado de máquina
            \end{itemize}
        \end{column}
        \begin{column}{0.48\textwidth}
            \textbf{Conexões:}
            \begin{itemize}
                \item Inferência Bayesiana
                \item Termodinâmica
                \item Mecânica estatística
                \item Princípio de mínima ação
            \end{itemize}
        \end{column}
    \end{columns}

    \vspace{0.3cm}
    \begin{center}
        \textit{Friston et al., Physics Reports, 2023}
    \end{center}
\end{frame}

\begin{frame}{IA para Resolver Equações Diferenciais}
    \begin{block}{Physics-Informed Neural Networks (PINNs)}
        \begin{equation*}
            \mathcal{L} = \mathcal{L}_{\text{dados}} + \lambda \mathcal{L}_{\text{física}}
        \end{equation*}

        Rede neural treinada para satisfazer EDPs
    \end{block}

    \begin{columns}[T]
        \begin{column}{0.48\textwidth}
            \textbf{Vantagens:}
            \begin{itemize}
                \item Sem malha (mesh-free)
                \item Alta dimensionalidade
                \item Dados escassos
                \item Problemas inversos
            \end{itemize}
        \end{column}
        \begin{column}{0.48\textwidth}
            \textbf{Aplicações:}
            \begin{itemize}
                \item Dinâmica de fluidos
                \item Mecânica quântica
                \item Propagação de ondas
                \item Relatividade geral
            \end{itemize}
        \end{column}
    \end{columns}

    \vspace{0.3cm}
    \begin{center}
        \textit{Karniadakis et al., Nature Reviews Physics, 2021}
    \end{center}
\end{frame}

\begin{frame}{IA para Descobrir Leis Físicas}
    \begin{block}{SINDy: Sparse Identification of Nonlinear Dynamics}
        \begin{equation*}
            \dot{\mathbf{x}} = \boldsymbol{\Theta}(\mathbf{x}) \, \boldsymbol{\xi}
        \end{equation*}
        onde $\boldsymbol{\Theta}$ é uma biblioteca de funções candidatas e $\boldsymbol{\xi}$ são coeficientes esparsos
    \end{block}

    \begin{columns}[T]
        \begin{column}{0.55\textwidth}
            \textbf{Exemplos de Sucesso:}
            \begin{itemize}
                \item Redescoberta de leis de Kepler
                \item Identificação de Hamiltonianos
                \item Leis de conservação
                \item Simetrias ocultas
            \end{itemize}
        \end{column}
        \begin{column}{0.42\textwidth}
            \begin{alertblock}{Questão Filosófica}
                \textit{A IA ``compreende'' as leis que descobre?}
            \end{alertblock}
        \end{column}
    \end{columns}

    \vspace{0.3cm}
    \begin{center}
        \textit{Brunton et al., PNAS, 2016}
    \end{center}
\end{frame}

\begin{frame}{Física Inspirando Arquiteturas de IA}
    \begin{columns}[T]
        \begin{column}{0.48\textwidth}
            \begin{block}{Lagrangian Neural Networks}
                \vspace{-0.2cm}
                \begin{equation*}
                    \frac{d}{dt}\frac{\partial \mathcal{L}_\theta}{\partial \dot{q}} = \frac{\partial \mathcal{L}_\theta}{\partial q}
                \end{equation*}
                \vspace{-0.4cm}
                \small Redes que respeitam estrutura lagrangiana
            \end{block}
            \vspace{-0.15cm}
            \begin{block}{Hamiltonian Neural Networks}
                \vspace{-0.2cm}
                \begin{equation*}
                    \dot{q} = \frac{\partial H_\theta}{\partial p}, \;\;
                    \dot{p} = -\frac{\partial H_\theta}{\partial q}
                \end{equation*}
            \end{block}
        \end{column}
        \begin{column}{0.48\textwidth}
            \textbf{Vantagens:}
            \begin{itemize}
                \item Conservação de energia garantida
                \item Estabilidade de longo prazo
                \item Extrapolação confiável
                \item Interpretabilidade física
            \end{itemize}

            \vspace{0.2cm}

            \textit{O melhor dos dois mundos: aprendizado com estrutura física}
        \end{column}
    \end{columns}
    \vspace{0.1cm}
    \begin{center}
        \small\textit{Cranmer et al., ICLR Workshop, 2020}
    \end{center}
\end{frame}

% ============================================
% SEÇÃO 5: IMPLICAÇÕES FILOSÓFICAS
% ============================================
\section{Implicações Filosóficas}

\begin{frame}{Complementaridade Epistemológica}
    \begin{center}
    \begin{tabular}{lcc}
        \toprule
        & \textbf{Mecânica Clássica} & \textbf{IA Generativa} \\
        \midrule
        Abordagem & Top-down & Bottom-up \\
        Raciocínio & Dedutivo & Indutivo \\
        Conhecimento & Leis fundamentais & Padrões nos dados \\
        Objetivo & Explicação causal & Predição estatística \\
        Transparência & Alta & Baixa \\
        Generalização & Por princípios & Por interpolação \\
        \bottomrule
    \end{tabular}
    \end{center}

    \vspace{0.5cm}

    \begin{alertblock}{Tese Central}
        \centering
        \textbf{Não são excludentes, mas complementares}
    \end{alertblock}
\end{frame}

\begin{frame}{Explicação vs. Predição}
    \begin{columns}[T]
        \begin{column}{0.48\textwidth}
            \begin{block}{Explicação Causal}
                \begin{itemize}
                    \item Por que o fenômeno ocorre?
                    \item Mecanismos subjacentes
                    \item Compreensão profunda
                    \item Transferível entre contextos
                \end{itemize}

                \vspace{0.3cm}
                \centering
                \textit{``A maçã cai porque a gravidade atrai''}
            \end{block}
        \end{column}
        \begin{column}{0.48\textwidth}
            \begin{block}{Predição Estatística}
                \begin{itemize}
                    \item O que vai acontecer?
                    \item Correlações observadas
                    \item Precisão prática
                    \item Contexto-dependente
                \end{itemize}

                \vspace{0.3cm}
                \centering
                \textit{``Historicamente, maçãs caem''}
            \end{block}
        \end{column}
    \end{columns}

    \vspace{0.5cm}

    \begin{center}
        \textbf{Ambas são formas válidas de conhecimento científico}
    \end{center}
\end{frame}

\begin{frame}{Hibridização: O Melhor dos Dois Mundos}
    \begin{center}
    \begin{tikzpicture}[scale=0.9, transform shape,
        box/.style={rectangle, draw=unespazul, fill=unespazul!20, thick,
                    minimum width=2.5cm, minimum height=1cm, align=center, rounded corners}]

        % Physics
        \node[box] (physics) at (-3,0) {Física\\(Estrutura)};

        % AI
        \node[box] (ai) at (3,0) {IA\\(Flexibilidade)};

        % Hybrid
        \node[box, fill=unespvermelho!20, draw=unespvermelho] (hybrid) at (0,-2) {Modelos\\Híbridos};

        % Arrows
        \draw[->, thick, unespazul] (physics) -- (hybrid);
        \draw[->, thick, unespazul] (ai) -- (hybrid);

        % Benefits
        \node[right=2cm of hybrid, align=left, font=\small] {
            $\bullet$ Conservação de energia\\
            $\bullet$ Interpretabilidade\\
            $\bullet$ Flexibilidade\\
            $\bullet$ Aprendizado eficiente
        };
    \end{tikzpicture}
    \end{center}

    \vspace{0.3cm}

    \begin{block}{Exemplos}
        \begin{itemize}
            \item Neural ODEs com estrutura hamiltoniana
            \item PINNs para problemas inversos
            \item Graph Neural Networks para simulações físicas
        \end{itemize}
    \end{block}
\end{frame}

\begin{frame}{IA Explicável Inspirada em Princípios Físicos}
    \begin{block}{Desafio da Interpretabilidade}
        Como tornar modelos de IA mais compreensíveis?
    \end{block}

    \begin{columns}[T]
        \begin{column}{0.48\textwidth}
            \textbf{Princípios da Física:}
            \begin{itemize}
                \item Leis de conservação
                \item Simetrias
                \item Localidade
                \item Causalidade
            \end{itemize}
        \end{column}
        \begin{column}{0.48\textwidth}
            \textbf{Aplicações em IA:}
            \begin{itemize}
                \item Equivariant NNs
                \item Conservation laws constraints
                \item Causal representation learning
                \item Physics-guided architectures
            \end{itemize}
        \end{column}
    \end{columns}

    \vspace{0.5cm}

    \begin{center}
        \textit{A física como fonte de princípios reguladores para IA}
    \end{center}
\end{frame}

\begin{frame}{Novos Paradigmas de Descoberta Científica}
    \begin{center}
        \Large\textbf{O Futuro da Física na Era da IA}
    \end{center}

    \vspace{0.5cm}

    \begin{enumerate}
        \item \textbf{Descoberta assistida por IA}
        \begin{itemize}
            \item Identificação de padrões em dados massivos
            \item Sugestão de hipóteses a serem testadas
        \end{itemize}

        \item \textbf{Simulação acelerada}
        \begin{itemize}
            \item Emuladores neurais de simulações físicas
            \item Exploração de espaços de parâmetros
        \end{itemize}

        \item \textbf{Síntese de conhecimento}
        \begin{itemize}
            \item Integração de dados heterogêneos
            \item Conexões entre áreas distintas
        \end{itemize}
    \end{enumerate}

    \vspace{0.3cm}

    \begin{alertblock}{}
        \centering
        \textit{O físico como integrador de paradigmas}
    \end{alertblock}
\end{frame}

% ============================================
% SEÇÃO 6: CONCLUSÃO
% ============================================
\section{Conclusão}

\begin{frame}{A Permanente Relevância da Mecânica Clássica}
    \begin{columns}[T]
        \begin{column}{0.55\textwidth}
            \begin{block}{Fundamento Conceitual}
                \begin{itemize}
                    \item Base de toda física moderna
                    \item Linguagem matemática universal
                    \item Paradigma de rigor científico
                    \item Framework para novos desenvolvimentos
                \end{itemize}
            \end{block}

            \begin{block}{Relevância Contemporânea}
                \begin{itemize}
                    \item Engenharia e tecnologia
                    \item Astrofísica e cosmologia
                    \item Biofísica e sistemas complexos
                    \item Inspiração para algoritmos
                \end{itemize}
            \end{block}
        \end{column}
        \begin{column}{0.42\textwidth}
            \begin{alertblock}{Mensagem Central}
                \textit{A Mecânica Clássica não é obsoleta --- é fundamental e atemporal}
            \end{alertblock}

            \vspace{0.5cm}

            \begin{quote}
                \small\itshape
                ``A mecânica de Newton é ainda a física que mais se usa.''
            \end{quote}
            \hfill --- Feynman
        \end{column}
    \end{columns}
\end{frame}

\begin{frame}{Mecânica Clássica na Era Digital}
    \begin{columns}[T]
        \begin{column}{0.48\textwidth}
            \begin{block}{Inspiração para IA}
                \begin{itemize}
                    \item Arquiteturas baseadas em física
                    \item Princípios variacionais
                    \item Leis de conservação
                    \item Simetrias como inductive bias
                \end{itemize}
            \end{block}
        \end{column}
        \begin{column}{0.48\textwidth}
            \begin{block}{Benchmark para IA}
                \begin{itemize}
                    \item Validação de modelos
                    \item Ground truth analítico
                    \item Teste de extrapolação
                    \item Avaliação de interpretabilidade
                \end{itemize}
            \end{block}
        \end{column}
    \end{columns}

    \vspace{0.5cm}

    \begin{center}
        \textbf{A Mecânica Clássica como pedra de toque para a era da IA}
    \end{center}
\end{frame}

\begin{frame}{Síntese Filosófica}
    \begin{center}
    \begin{tikzpicture}[scale=0.85, transform shape]
        % Venn diagram
        \draw[thick, fill=unespazul!30] (-1.5,0) circle (2.5cm);
        \draw[thick, fill=unespvermelho!30] (1.5,0) circle (2.5cm);

        % Labels
        \node at (-3,2.5) {\textbf{Mecânica Clássica}};
        \node at (3,2.5) {\textbf{IA Generativa}};

        % Left circle content
        \node[align=center] at (-2.8,0) {\small Determinismo\\\small Causalidade\\\small Leis universais};

        % Right circle content
        \node[align=center] at (2.8,0) {\small Probabilismo\\\small Correlação\\\small Padrões locais};

        % Intersection
        \node[align=center] at (0,0) {\small Otimização\\\small Modelagem\\\small Predição};

        % Bottom message
        \node at (0,-3.5) {\large\textbf{Complementaridade Epistemológica}};
    \end{tikzpicture}
    \end{center}
\end{frame}

\begin{frame}{Mensagem Final}
    \begin{center}
        \Large
        \textbf{Diferentes formas de conhecimento não são excludentes}

        \vspace{0.6cm}

        \textbf{A ciência avança pela pluralidade metodológica}

        \vspace{0.6cm}

        \textbf{O físico contemporâneo: integrador de paradigmas}
    \end{center}

    \vspace{0.5cm}

    \begin{alertblock}{}
        \centering
        A Mecânica Clássica permanece relevante não \textit{apesar} da IA, mas \textit{junto} com ela --- como instrumento complementar na busca pela compreensão da natureza.
    \end{alertblock}
\end{frame}

\begin{frame}
    \centering
    \vfill

    {\Huge\textbf{Obrigado}}

    \vspace{1cm}

    {\Large Perguntas?}

    \vfill

    \begin{quote}
        \itshape
        ``O objetivo da física é descobrir as leis que a natureza escolheu, não as leis que nós gostaríamos que ela tivesse escolhido.''
    \end{quote}
    \hfill --- Richard Feynman

    \vfill
\end{frame}

% ============================================
% REFERÊNCIAS
% ============================================
\begin{frame}[allowframebreaks]{Referências Selecionadas}
    \footnotesize

    \textbf{Mecânica Clássica}
    \begin{itemize}
        \item GOLDSTEIN, H. \textit{Classical Mechanics}. 2nd ed. Addison Wesley, 2002.
        \item LANDAU, L. D.; LIFSHITZ, E. M. \textit{Mechanics}. 3rd ed. Butterworth-Heinemann, 1976.
        \item ARNOLD, V. I. \textit{Mathematical Methods of Classical Mechanics}. Springer, 1989.
    \end{itemize}

    \framebreak

    \textbf{IA e Aprendizado de Máquina}
    \begin{itemize}
        \item GOODFELLOW, I.; BENGIO, Y.; COURVILLE, A. \textit{Deep Learning}. MIT Press, 2016.
        \item VASWANI, A. et al. Attention Is All You Need. \textit{NeurIPS}, 2017.
    \end{itemize}

    \textbf{Conexões Física-IA}
    \begin{itemize}
        \item KARNIADAKIS, G. E. et al. Physics-informed machine learning. \textit{Nature Reviews Physics}, v. 3, p. 422-440, 2021.
        \item BRUNTON, S. L. et al. Discovering governing equations from data. \textit{PNAS}, v. 113, n. 15, p. 3932-3937, 2016.
        \item FRISTON, K. et al. The free energy principle made simpler. \textit{Physics Reports}, v. 1024, p. 1-29, 2023.
        \item CRANMER, M. D. et al. Lagrangian Neural Networks. \textit{ICLR Workshop}, 2020.
    \end{itemize}
\end{frame}

\end{document}
